\section{논의}

본 실험의 질문은 단순히 ``하이브리드 로봇이 걷고 주행할 수 있는가''가 아니다.
하나의 구동 말단 접촉부를 재사용해 어떤 거동을 얻을 수 있는지를 묻는다. 현재
개방형 호 모델에서 말단 단독 회전은 $150\,$mm 계단을 통과했고, 시험한 제어기
중에서는 협응 관절 운동만 $200\,$mm 계단을 통과했다. 평지에서는 말단 회전만
사용한 조건이 관절 보행보다 느리고 기계적 비용도 컸으며, 몸체 자세와 말단 회전을
함께 명령한 조건의 기술적 평균 속도가 가장 높았다. 이 결과는 하나의 링크를
지지·준구름·모서리 접촉 후보로 재사용한다는 가설을 뒷받침하지만, 계단 결과를
개구부에만 귀속하지는 못한다. 식~(\ref{eq:stairmoment-ko})이 보여주듯 축토크로
구동되는 바퀴에는 수평력만 사용하는 경우의 특이점이 적용되지 않는다. 따라서
개방형 호와 닫힌 바퀴를 동일 조건에서 비교하는 실험이 결정적인 검증이다.

비용도 명확하다. 협응 호 이동은 관절 보행보다 미끄럼이 $6.23\times$, 기계일이
$3.18\times$ 컸다. 식~(\ref{eq:contactspeed-ko})은 그 1차 원인을 설명한다.
현재 제어기는 접촉 속도 피드백 대신 명령 활성도에 비례하는 표면 속도를 사용하므로,
명령한 표면 속도가 몸체 속도를 크게 초과할 수 있다. 따라서 현재 모델은 속도와
비용의 절충을 보여줄 뿐 에너지 효율이 높다는 주장을 뒷받침하지 않는다. 몸체 및
접촉 속도에 연동한 회전 제어는 형상 절제와 분리해 검증해야 할 제어기 개선안이다.

결과를 일반화할 때에는 네 가지 한계를 고려해야 한다. 첫째, 모든 증거는 결정론적
시뮬레이션에서 얻었다. TPU의 접촉 연성과 마찰, 폴리머 채움률에 따른 질량, 기어
효율, 백래시, 열적 한계, 배터리 전압 강하는 근사값이며, 모델 총질량도 재조립 후
직접 계량해야 한다. 둘째, 공칭 조건은 조건별 1회이므로 신뢰구간이 없다. 셋째,
보행과 협응 호 이동의 주파수와 스트로크 한계가 서로 달라 동일 파라미터 비교가
아니다. 넷째, $150\,$mm 결과의 모서리 걸림 해석은 기하 관계와 관절 운동을 쓰지
않았다는 사실에 근거한 추론이지 직접 형상 절제가 아니다. \ref{sec:method-ko}절의
고정 짝비교 프로토콜은 시뮬레이션에서 뒤의 세 한계를 다루지만, 실물 모델 식별은
여전히 필요하다. 계단 상위 제어기 또한 계단 형상을 미리 입력받으며 지형 인식이나
발판 탐색을 수행하지 않는다.

\section{결론}

본 논문은 여섯 개의 개방형 호형 말단을 지지 발, 준구름 접촉부, 계단 모서리 접촉
후보로 재사용하는 18액추에이터 6족 로봇과, 삼각보행과 연속 호 회전을 같은 말단
관절에 동시에 명령하는 제어기를 제시했다. 결정론적 MuJoCo 단일 시행에서 관절
보행의 평균 속도는 $0.117\,\mathrm{m\,s^{-1}}$, 협응 제어기는
$0.200\,\mathrm{m\,s^{-1}}$였지만 협응 조건의 기계일과 미끄럼이 더 컸다.
개방형 호 계단 시험에서는 관절 운동 없이 $150\,$mm 형상을 통과했고, $200\,$mm
형상은 협응 관절 운동을 사용한 조건만 통과했다. 전환 조건은 직립을 유지하며
$0.7\,$s 이내에 완료되었다. 이 단일 시행 결과는 가설과 재현 가능한 운용 범위를
제시할 뿐 형상의 인과 효과를 입증하지 않는다. 함께 제공하는 동일 형상 조건 및
짝비교 프로토콜은 남은 개방형 호 주장을 반증 가능한 형태로 만든다. 실물 속도,
내구성, 에너지 효율, 계단 등반 신뢰도는 아직 검증하지 않았다.
